\section{行在不是临时的空名}

1127年赵构即位，高宗朝廷开始重建。其后的驻跸扬州、组织勤王与财政，1129年苗傅、刘正彦兵变，1130年金军渡江后北撤，说明南渡首先是一个如何让诏令、军队、仓储与税源重新接合的问题。南宋本纪和《建炎以来系年要录》为1127--1162年提供密集的宋方编年脊梁；它们所记录的战报与奏议，不能自动说明地方社会或金军的全部经验。

1132年行在逐渐以临安为稳定基地，1138年临安被确立为行在和事实上的都城。这里的“稳定”并不意味着南方不再流动，而是朝廷能将机构、市场、漕运和军事调度更持续地集中在江南。临安的城市空间、供给和人口变化，须由遗址、契约、地方文献与制度记录交叉考察；不能仅根据京师称号，把它写成对开封的无缝替代。

1134年宋军在襄汉、淮西有局部胜利，1135年岳飞等将领获得更大经略与整编责任，1140年又发动反攻。局部战果并未消除中枢对兵权、粮饷与和议的顾虑。1141年南宋接受绍兴和议框架，1142年岳飞被处死、相关将领和军队被整顿。人物传记与后世纪念常把这一组事件压成单一忠奸故事；《宋史》、要录和会要更适合用来分别追问和约、任免、军队整编和法制处置，不能据史家褒贬补出未见于材料的宫廷动机。

\section{财政、和议与再一次北伐}

1150年朝廷调整会子和纸币管理，提醒我们战争后的国家不只靠边军维持。会子制度的颁行、改限和回收可以由会要与食货材料追踪，发行额却不等于流通量，亦不等于不同地区家庭承担的物价风险。1161年金军渡江受挫，1162年高宗禅位、孝宗即位；1163年孝宗支持北伐而在符离等地受挫，1164年隆兴和议订立。继承、战争和货币管理彼此相连，但没有任何一项可单独解释另两项。

1173年朱熹等在书院、讲学与地方事务中活动，1194年宁宗即位后韩侂胄在中枢取得影响，1196年庆元党禁限制朱熹及相关士人的政治和学术活动。书院与讲学不能化约为脱离政务的“文化”，党禁也不能由一张名单推出士人社会的一致反应。保存较多的是精英文本和朝廷处分；地方学校、出版和社群如何回应，须由区域材料分别回答。

1204年北伐准备开始，1206年开禧北伐不利，1207年韩侂胄被杀，1208年嘉定和议结束战事。北伐的失败并非只由一人或一战决定，而牵涉动员、前线协作、金的反应和中枢政治。1217年金军南攻四川、两淮，表明和议之下的边界仍须投入军费与运输。南宋的国家能力正在这种反复中形成：它能在江南维持税源和行政，却不能把北方议程变成低成本的选择。

\section{从攻金联盟到蒙古前线}

1224年理宗即位、史弥远继续主导中枢；1233年南宋与蒙古就攻金短暂协作，1234年蔡州陷落后宋军尝试进入河南故地。将协作视作两国拥有同一长期目标，会遮蔽金亡后立即出现的道路、城池与粮区竞争。1235年蒙古开始多方向攻宋，1253年大理被征服而西南战略环境改变，1258年蒙哥发动大规模攻宋计划，1259年钓鱼城战事与蒙哥去世改变战争节奏。宋、蒙叙事对这些节点的解释不同，战场伤亡和个人行动尤其不可未经互校。

1264年度宗即位、贾似道权势上升；1268年襄阳、樊城受围，1273年失守，1274年元军大举南下。襄樊不是孤立的城战，而是汉水、川鄂道路、江南财政与海陆运输之间的关节。1275年宋军在丁家洲等失利、贾似道被罢，1276年临安降元；遗臣、宗室和地方军民仍在东南抵抗，直至1279年崖山后南宋政权终结。降城、抵抗与海战的叙述都带有强烈的正统和忠义语言；应分开朝廷的终结、地方控制的转换与不同人群的去留，而不是把它们归为同一时刻。
